\chapter{Setup and preparatory results}\label{chap:setup}

\section{Notation}

Let $n\ge 2$ and $d=\binom{n}{2}$.
An $n\times 4t$ partial Hadamard matrix has $\pm1$ entries with
pairwise orthogonal rows.
Let $N_{n,4t}$ be the number of such matrices and
$A_{n,4t}=\binom{2n}{n}^{2t}$ the Gaussian approximation.
Set $\psi(\lambda)=\E[\exp(i\sum_{i<j}\lambda_{\{i,j\}}\xi_i\xi_j)]$
for independent Rademacher signs $\xi_i$,
$s(\lambda)=\|\lambda\|^2=\sum_e\lambda_e^2$, and
\[
  K_n:=2^{2d-n+1}(2\pi)^{-d},
  \qquad
  \core(d,t):=\int_{\corereg}e^{-2t\|\lambda\|^2}\,d\lambda,
  \qquad
  F(d,t):=\left(\frac{\pi}{2t}\right)^{d/2}.
\]

\section{Results from de~Launey--Levin}

\begin{lemma}[Cosine-product bound]\label{fact:psi-sq}
\lean{universal_magnitude_bound_full}\leanok
For every $\gamma\in[-\pi,\pi]^d$ and every $k\in\{1,\dots,n\}$,
$|\psi(\gamma)|^2\le \frac12+\frac12\prod_{i\ne k}\cos(2\gamma_{\{i,k\}})$.
\end{lemma}

\begin{lemma}[Lattice structure]\label{fact:lambda-facts}
\leanok
The formalization proves the following (as private lemmas in
\texttt{HadamardCn3TorusCount.lean}):
\begin{enumerate}
\item \emph{Multiplicativity:}
$\psi(\lambda+\gamma)=\psi(\lambda)\psi(\gamma)$ for $\lambda\in\Lambda$,
$\gamma\in\Bbox_{\pi/4}$
\textup{(\texttt{psi\_translate\_lambdaReal})}.
\item \emph{Disjoint boxes:}
the translates $\Bbox_\delta(\lambda)$, $\lambda\in\Lambda$, are
pairwise disjoint for $\delta<\pi/4$
\textup{(\texttt{translated\_edgeBoxes\_disjoint})}.
\item \emph{Cardinality:}
$|\Lambda|=2^{2d-n+1}$ and $\psi(\lambda)^{4t}=1$ for $\lambda\in\Lambda$
\textup{(\texttt{card\_lambdaShifts\_eq\_pow\_of\_two\_le},
\texttt{exp\_phase\_lambdaReal\_false\_pow\_four\_mul})}.
\item \emph{Even-degree characterization:}
$\lambda\in\Lambda$ iff every vertex of $G_{\lambda^{(2)}}$ has
even degree
\textup{(\texttt{parityEvenBool\_iff\_incidence\_zero})}.
\end{enumerate}
These are private because they are only used internally
by the torus counting machinery.
\end{lemma}

\begin{proposition}[Primary-secondary decomposition]\label{prop:primary-secondary}
\lean{normalizedCount_primary_secondary_decomposition}\leanok
\uses{fact:lambda-facts}
For every $t\ge 1$ and every $0<\delta<\pi/4$,
\[
  \Prb(S_{4t}=0)
  =
  2^{2d-n+1}(2\pi)^{-d}\int_{\Bbox_\delta}\psi(\lambda)^{4t}\,d\lambda
  +(2\pi)^{-d}\int_{R_\delta}\psi(\gamma)^{4t}\,d\gamma.
\]
\end{proposition}

\begin{lemma}[Residual decomposition]\label{lem:residual-decomp}
\lean{normalizedCount_primary_secondary_decomposition}\leanok
\uses{fact:lambda-facts}
For $0<\delta<\pi/4$,
\[
  R_\delta = R_\delta^{\mathrm{even}} \cup R_\delta^{\mathrm{odd}},
  \qquad
  R_\delta^{\mathrm{even}}
  = \bigcup_{\lambda\in\Lambda}\bigl(\lambda+[\Bbox_{\pi/4}\setminus\Bbox_\delta]\bigr),
  \qquad
  R_\delta^{\mathrm{odd}}
  = R_\delta\setminus R_\delta^{\mathrm{even}},
\]
and the sets in the union are disjoint.
\end{lemma}

\begin{proof}\leanok
\uses{fact:lambda-facts}
The translates $\Bbox_{\pi/4}(\lambda)$, $\lambda\in\Lambda_0$, tile $\TT^d$.
Even cells centered at $\lambda\in\Lambda$ contribute to $R_\delta^{\mathrm{even}}$;
odd cells centered at $\lambda\in\Lambda_0\setminus\Lambda$ contribute to
$R_\delta^{\mathrm{odd}}$.
\end{proof}

\begin{lemma}[Odd-cell bound]\label{lem:odd-bound}
\lean{universal_magnitude_bound_full}\leanok
\uses{fact:psi-sq,fact:lambda-facts,lem:residual-decomp}
If $\gamma\in R_\delta^{\mathrm{odd}}$, then $|\psi(\gamma)|^2\le \tfrac12$.
\end{lemma}

\begin{proof}\leanok
\uses{fact:psi-sq,fact:lambda-facts}
Since $\gamma$ lies in a quarter-box centered at $\lambda\in\Lambda_0\setminus\Lambda$,
some vertex $k$ has odd degree in $G_{\lambda^{(2)}}$,
so $\prod_{i\ne k}\cos(2\gamma_{\{i,k\}})\le 0$,
and the cosine-product bound (Lemma~\ref{fact:psi-sq}) gives
$|\psi(\gamma)|^2\le 1/2$.
\end{proof}

\begin{lemma}[Positivity of $\RePart\psi$]\label{lem:realpart}
\lean{re_psi_taylor4_decomposition}\leanok
If $s(\lambda)\le 1/2$, then $\RePart\psi(\lambda)\in[3/4,1]$.
\end{lemma}

\begin{proof}\leanok
$\RePart\psi(\lambda)=\E[\cos X_\lambda]\ge 1-\frac12 s(\lambda)\ge 3/4$.
\end{proof}

\section{Fixed-$n$ asymptotics}

\begin{lemma}[Fixed-$n$ asymptotics]\label{fact:fixed-n}
\lean{fixed_n_count_asymptotic}\leanok
For every fixed integer $n\ge 2$,
$N_{n,4t}=[1+o_{t\to\infty}(1)]\,A_{n,4t}$ as $t\to\infty$.
\end{lemma}

\begin{proof}\leanok
Shrinking-box argument with $\delta_t=t^{-2/5}$: the primary integral
over $\Bbox_{\delta_t}$ converges to $F(d,t)$ as $t\to\infty$,
and the residual is negligible for any fixed $n$.
\end{proof}

\section{Standard analytic tools}

\begin{lemma}[Hypercontractive inequality]\label{lem:hc}
\lean{fixedDegreeHC_degree2_W_fourth}\leanok
Let $H$ be a polynomial of degree at most $q$ in either independent
Rademacher signs or independent centered Gaussians.
Then for every $p\ge 2$,
$\|H\|_{L^p}\le (p-1)^{q/2}\|H\|_{L^2}$.
\end{lemma}

\begin{lemma}[Gaussian radial moments]\label{lem:gaussian-radial}
\lean{gaussian_radial_moments}\leanok
For every integer $m\ge 0$ there is $C_m>0$ such that
$\int_{\R^d}\|\lambda\|^{2m}e^{-2t\|\lambda\|^2}\,d\lambda
\le C_m(d/t)^m F(d,t)$ for all $d,t\ge 1$.
\end{lemma}

\begin{lemma}[Gaussian quadratic integral]\label{lem:gaussian-quadratic}
\lean{gaussian_integral_formula}\leanok
Let $A$ be a complex symmetric $m\times m$ matrix whose real part is
positive definite. Then
$\int_{\R^m}e^{-x^\top A x}\,dx=\pi^{m/2}\det(A)^{-1/2}$.
In particular, if $g\sim N(0,I_m)$ and $M$ is real symmetric,
$|\E[e^{ig^\top Mg/2}]|=\det(I+M^2)^{-1/4}$.
\end{lemma}

\begin{proof}\leanok
Diagonalize $S=\RePart(A)^{-1/2}\,\mathrm{Im}(A)\,\RePart(A)^{-1/2}$,
factor the integral into one-dimensional Fresnel integrals,
and reassemble via $\det(A)=\det(\RePart A)\det(I+iS)$.
\end{proof}

\begin{lemma}[Core Gaussian mass]\label{lem:inner-core}
\lean{inner_core_gaussian_mass}\leanok
There exists an absolute constant $c_*>0$ such that
$\core(d,t)\ge (1-e^{-c_*d})F(d,t)$.
\end{lemma}

\begin{proof}\leanok
On $\|\lambda\|^2>d/t$, one has $e^{-2t\|\lambda\|^2}\le e^{-d}e^{-t\|\lambda\|^2}$,
so the tail is at most $e^{-d}(\pi/t)^{d/2}\le e^{-c_*d}F(d,t)$
with $c_*=1-\frac12\log 2$.
\end{proof}
